\cvsection{Expériences professionnelles}

\begin{cventries}
  \cventry
    {Stagiaire ingénieur de recherche en IA}
    {INRIA}
    {Bordeaux, France}
    {Mars 2026 -- Septembre 2026}
    {
      \begin{cvitems}
        \item {Conception en autonomie, avec le suivi de mes encadrants, d’un agent IA multi-étapes répondant au besoin d’aider les chercheurs à trouver des articles scientifiques.}
        \item {Construction de l’agent sans framework d’orchestration afin de maîtriser ses outils, son état, sa mémoire, son contexte et ses enchaînements de tâches.}
        \item {Développement d’un mécanisme permettant à l’utilisateur de sélectionner les messages à inclure dans le contexte transmis au LLM afin de contrôler la mémoire de la conversation.}
        \item {Intégration de l’API OpenAI et conception de prompts pour adapter le comportement de l’agent et guider la sélection de ses outils.}
        \item {Collecte automatisée de documents scientifiques en ligne, préparation et découpage en segments adaptés au pipeline RAG.}
        \item {Génération des embeddings, indexation dans ChromaDB et recherche sémantique des passages pertinents pour enrichir le contexte du LLM.}
        \item {Développement du backend Python/FastAPI avec API REST, OpenAPI et MongoDB, puis d’une interface React/Next.js destinée aux utilisateurs.}
        \item {Création de webhooks transmettant les résultats de tâches cron et intégration des différents outils à l’application.}
        \item {Installation de modèles locaux et utilisation de vLLM pour les lancer et les servir sur des GPU NVIDIA.}
        \item {Comparaison de Qwen, Gemma et GLM et de plusieurs quantifications selon la qualité, la latence, le TTFT et la taille du contexte.}
        \item {Conteneurisation de composants avec Docker et mise en place d’un pipeline CI/CD sur GitHub.}
      \end{cvitems}
    }

  \cventry
    {Stagiaire de recherche en apprentissage par renforcement}
    {Université Karunya}
    {Coimbatore, Tamil Nadu, Inde}
    {Juin 2025 -- Septembre 2025}
    {
      \begin{cvitems}
        \item {Développement d’un modèle d’apprentissage par renforcement et de son environnement d’entraînement pour la navigation autonome d’un robot hospitalier.}
        \item {Déploiement et optimisation du logiciel et du modèle sur un Raspberry Pi embarqué.}
      \end{cvitems}
    }
\end{cventries}
